\chapter{Chapter Trial}
\label{chp-trial}

\section{Purpose}

	\begin{enumerate}
		\item To judge infractions of the bylaws on a chapter-wide level.
	\end{enumerate}

\section{Membership}
	
	\begin{enumerate}
		\item The President of the chapter presides over the trial. If the President is unable to serve or not present, the Arbiter will preside over the trial.
		\item If neither the President nor the Arbiter are able to serve, follow the Order of Succession of Standards Board \hyperref[stds-board-succession]{Article XII, Section D, \autoref*{stds-board-succession}})
	\end{enumerate}

\section{Voting}
	
	\begin{enumerate}
		\item Unless otherwise specified, votes during a Chapter Trial are by simple majority.
		\item Every member of the chapter who is present for the trial except for the Scribe and Parliamentarian is a voting member.
		\item The Scribe and Parliamentarian are non-voting members, and the member who is presiding over the trial will only vote in the case of a tie.
		\item Quorum of the trial consists of $\frac{2}{3}$ of the chapter.
	\end{enumerate}

\section{Charges}
	
	\begin{enumerate}
		\item When the accusation is received by the Arbiter, Scribe, and Parliamentarian, they will collectively decide whether the issue will be brought to standards, or the chapter.
			\begin{enumerate}
				\item If the accuser and/or accused disagree with the decision of which trial the charges are brought to, they may appeal the decision.
				\item The appeal will be brought exclusively to the standards board justices to vote on whether the appeal is valid.
			\end{enumerate}

	\end{enumerate}

\section{Accusation Procedure}

	\begin{enumerate}
		\item All accusations must be submitted in writing to the Arbiter wihtin one month of the most recent violation.
			\begin{enumerate}
				\item Each accusation must be as detailed as possible including at least the clause of the local bylaws that was violated.
				\item Each accusation must be submitted within one month of the violation.
			\end{enumerate}

		\item Upon receiving the Accusation, the Arbiter, Scribe, and Parliamentarian will determine the validity of the claim.
		\item The Arbiter will inform the accuser of the validity of the accusation.
		\item If the accusation is found to be without merit, proceedings will end and no hearing will be held.
		\item If the accusation is found to have merit, the Arbiter will provide written notification to the accused.
			\begin{enumerate}
				\item The notification must be given within forty-eight hours of receipt of the accusation.
				\item The notification must include the charges with enough detail to allow the accused to prepare a defense.
			\end{enumerate}
		\item Upon receiving the notification, the accused has twenty-four hours to submit a plea in writing to the Arbiter.
			\begin{enumerate}
				\item The possible pleas are \gls{res-plea}, \gls{not-res-plea}, and \gls{no-cont-plea}.
				\item If a plea is not received, a \gls{no-cont-plea} will be entered.
			\end{enumerate}

		\item After a plea is entered, the Arbiter will schedule a time and place for the hearing.
			\begin{enumerate}
				\item The hearing shall occur between 1 and 3 weeks after a plea is entered.
				\item The Arbiter will notify the accused, accuser, and the board of the time, place, and nature of the hearing. The accusation an dplea will also be submitted to the Standards Board Justices.
			\end{enumerate}

		\item The Parliamentarian will inform \gls{ec} of the accusation, the accused, and the nature of the complaint at the next \gls{ec} meeting.

		% Should this section still be included?
		\item Recusal
			\begin{enumerate}
				\item Members of the board will recuse themselves from the hearing if they feel a conflict of interest exists or they are the accused.
				\item Members of the board may be recused by \gls{simp-maj} vote of the remaining board. The accused, accuser, or a member of the board may call for such a vote.
				\item Justices and members of the Standards Board may recuse themselves or others up to the beginning of the trial.
			\end{enumerate}

	\end{enumerate}

\section{Hearing Preparation}

	\begin{enumerate}
		\item Evidence will be distributed by the Arbiter between twenty-four and forty-eight hours before the trial date and time.
		\item Attire will be at the discretion of the Arbiter.
		\item Witnesses may volunteer to speak on behalf of or against the charges.
			\begin{enumeration}
				\item The witnesses may volunteer by giving a written message to the Arbiter stating their intent to participate and stance at least three hours before the trial.
				\item Witnesses will be advised by the Arbiter that they are not to discuss the hearing or the accusation outside of the hearing.
				\item The Arbiter will call the hearing to order.
			\end{enumeration}

	\end{enumerate}

\section{Hearing Procedure}

	\begin{enumerate}
		\item The Arbiter reads aloud the charges and plea and ensures that the accused understands the charges and their plea.
		\item The accuser may give an opening statement with a time limit as voted on by the chapter.
		\item The accused may give an opening statement with a time limit as voted on by the chapter.
		\item The accuser may then be asked questions by the chapter with a time limit as voted on by the chapter.
		\item The accused may then be asked questions by the chapter with a time limit as voted on by the chapter.
		\item All witnesses will be brought into the room by the Arbiter.
		\item Witnesses that wish to speak for the charges will be chosen in an order decided by the Arbiter to give a statement. The accused and the justices will ask questions of the witness(es) with a time limit as voted on by the chapter.
		\item Witnesses that wish to speak against the charges will be chosen in an order decided by the Arbiter to give a statement. The accused and the justices will ask questions of the witness(es) with a time limit as voted on by the chapter.
		\item All witnesses will be excused by the Arbiter.
		\item The accuser may give a closing statement with a time limit as voted on by the chapter.
		\item The accused may give a closing statement with a time limit as voted on by the chapter.
		\item The Arbiter will excuse the accuser and the accused from the hearing.
		\item The chapter will then determine whether the accused is responsible or not responsible.
		\item The chapter will then determine appropriate sanctions, if any.
			\begin{enumeration}
				\item Past judicial history of the accused and sanctions given for similar accusations in the past may be considered.
				\item There will be a discussion to determine each specific sanction. After discussion comes to a close, the sanctions decided upon must be passed by a $\frac{2}{3}$ of present members vote. If the vote is not passed, the chapter will return to discussion.
			\end{enumeration}
		\item After the verdict and sanctions have been determined the Arbiter will readmit the accused and accuser back to the hearing and will announce the verdict and sanctions given, if any.
		\item The Arbiter will then inform the accused of the appeal process.
		\item The Arbiter will then adjourn the hearing.
	\end{enumerate}	

\section{Sanctions}
	
	\begin{enumerate}
		\item The chapter may issue sanctions within its scope. The following are suggested examples of possible sanctions:
			\begin{enumeration}
				\item Counseling
				\item Loss of housing status
				\item Loss of voting privileges
				\item Recommendation of Removal from an Elected Position
				\item Recommendation of Suspension
				\item Probation
			\end{enumeration}

	\end{enumerate}